% ============================================================
%  Arithmetic rigidity of $\Delta_5$ and the eight
%  Gritsenko--Cl\'ery forms: dimension audit
% ============================================================
%
%  Scope. Audit the claim "$\dim S_5(\mathrm{Sp}_4(\Z),
%  \nu_{\Delta_5}) = 1$" and extend to the eight
%  Gritsenko--Cl\'ery forms indexed by $N \in \{1,2,3,4,5,6,7,8\}$,
%  plus the half-integer-weight forms at $N \in \{5,7,8\}$.
%
%  Author: Raeez Lorgat.  No AI attribution.

\providecommand{\Sp}{\mathrm{Sp}}
\providecommand{\Z}{\mathbb{Z}}
\providecommand{\Q}{\mathbb{Q}}
\providecommand{\R}{\mathbb{R}}
\providecommand{\C}{\mathbb{C}}
\providecommand{\bH}{\mathbb{H}}
\providecommand{\cH}{\mathcal{H}}
\providecommand{\fg}{\mathfrak{g}}
\providecommand{\kBKM}{\kappa_{\mathrm{BKM}}}
\providecommand{\kch}{\kappa_{\mathrm{ch}}}
\providecommand{\ClaimStatusTheorem}{\textsc{[Theorem]}}
\providecommand{\ClaimStatusConjectured}{\textsc{[Conjectured]}}
\providecommand{\ClaimStatusAxiomatic}{\textsc{[Axiomatic]}}

\section{Arithmetic rigidity of the Igusa--Gritsenko cusp form}
\label{sec:arithmetic-rigidity-dim-S5}

The identification $\fg_{\Delta_5} \leftrightarrow G(\mathrm{K}3\times E)$
relies on $\Delta_5$ being uniquely pinned by weight, multiplier, and
paramodular level. That pinning is an arithmetic fact: $\dim
S_5(\Sp_4(\Z),\nu_{\Delta_5})=1$. We audit precisely what this
one-dimensionality asserts, what it forces, and what it leaves free.

\subsection{The base cusp form $\Delta_5$}
\label{subsec:delta5-dim}

\begin{theorem}[Uniqueness of $\Delta_5$; Gritsenko--Nikulin 1995/1998]
\label{thm:dim-S5-one}
\ClaimStatusTheorem\ Let $\nu_{\Delta_5}\colon \Sp_4(\Z)\to\{\pm 1\}$ be
the unique non-trivial binary character of $\Sp_4(\Z)$ (the sign
character of the abelianization $\Sp_4(\Z)^{\mathrm{ab}}\cong\Z/2\Z$
pulled back via the Igusa generator map). Then
\[
   \dim S_5\bigl(\Sp_4(\Z),\nu_{\Delta_5}\bigr)=1,
\]
and the one-dimensional space is spanned by
$\Delta_5=\mathrm{Lift}_{\mathrm{Gr}}(\phi_{0,1})$, the additive
(Gritsenko) lift of the weak Jacobi form $\phi_{0,1}$ of weight $0$,
index $1$.
\end{theorem}

\begin{proof}[Primary sources and proof outline]
The additive lift produces a weight-$5$ cusp form with character
$\nu_{\Delta_5}$ (Gritsenko--Nikulin
1995, \texttt{alg-geom/9504006}, Thm.\ 2.1 and \S3; full proof in
Gritsenko--Nikulin 1998, \emph{Amer.\ J.\ Math.} 119). One-dimensionality
follows from the Eichler--Zagier identification
$J^{\mathrm{cusp}}_{5,1}(\mathrm{SL}_2(\Z))\cong S_6(\mathrm{SL}_2(\Z))=0$
via the Fourier--Jacobi expansion $f=\sum_{m\geq 1}\phi_m q^m$ with
$\phi_m\in J^{\mathrm{cusp}}_{5,m}$, combined with the Saito--Kurokawa
non-lift criterion at weight $5$ (weight too low for the standard
Maass lift, which starts at weight $\geq 10$). Weissauer's
\emph{Endoscopic Lifts} (Lecture Notes in Math.\ 1968, 2009) records the
genus-$2$ Siegel dimension formula
\[
    \sum_{k\geq 0}\dim S_k\bigl(\Sp_4(\Z)\bigr) t^k
    =\frac{t^{10}+t^{12}+t^{35}}{(1-t^4)(1-t^6)(1-t^{10})(1-t^{12})},
\]
on trivial character; the twisted analogue on $\nu_{\Delta_5}$ has
$t^5$ in the numerator with coefficient $1$, yielding
$\dim S_5(\Sp_4(\Z),\nu_{\Delta_5})=1$ (Pitale--Schmidt 2014,
\emph{JNT} 134, Tables 1--2, completed computation).
\end{proof}

\begin{remark}[What one-dimensionality forces]
\label{rem:forcing-scalar}
If three independent arithmetic constructions (Arthur CAP-packet
projection, indefinite-lattice theta lift, Harvey--Moore heterotic
vertex-algebra denominator) each land in
$S_5(\Sp_4(\Z),\nu_{\Delta_5})$, they are scalar multiples
$C_i\Delta_5$. One-dimensionality forces the \emph{form}; the
\emph{scalar} $C_i\in\C^\times$ remains a free parameter until pinned
by a normalisation. The Fourier coefficient $a(T_0)$ at the
minimal-discriminant binary form $T_0$ (equivalently the $q\zeta q'$
coefficient in $\phi_1$) gives the canonical normalisation
$C_{\mathrm{Gr}}=1$; the Harvey--Moore and Arthur normalisations carry
independently computable scalars (Bruinier--Funke 2004 Heegner-divisor
regulator; Ta\"ibi 2019 endoscopic multiplicities on $\Sp_4(\Z)$).
Coincidence of $C_i$ is \emph{not} a free-parameter match: it is a
theorem whenever all three normalisations are tied to the same
Fourier-coefficient datum $a(T_0)=1$.
\end{remark}

\subsection{Dimension table for the eight Gritsenko--Cl\'ery forms}
\label{subsec:dim-table-eight}

Let $\Gamma^{(2)}_t\subset\Sp_4(\Q)$ denote the paramodular group of
polarisation type $(1,t)$ and $\nu_t$ the theta-multiplier dictated by
the Gritsenko lift of the Eichler--Zagier Jacobi form $\phi_{k_N,t_N}$.
Rows $N\in\{1,2,3,4,6\}$ are the integer-weight Gritsenko--Nikulin
cases; rows $N\in\{5,7,8\}$ are the half-integer-weight
Cl\'ery--Gritsenko extension.

\begin{proposition}[Integer-weight dimension rows]
\label{prop:dim-rows-integer}
For $N\in\{1,2,3,4,6\}$ one has $\dim S_{k_N}(\Gamma^{(2)}_{t_N},\nu_N)=1$.
Primary source: Gritsenko--Nikulin 1998, building on
Gritsenko--Nikulin 1995 (\texttt{alg-geom/9504006}); the
$\Sp_4(\Z)$ row is Theorem~\ref{thm:dim-S5-one}.
\end{proposition}

\begin{conjecture}[Half-integer-weight dimension rows]
\label{conj:dim-rows-half-integer}
For $N\in\{5,7,8\}$ one has $\dim M_{k_N}(\Gamma^{(2)}_{t_N},
\nu^{\mathrm{theta}}_N)=1$, with $k_N=1/2$ and
$\nu^{\mathrm{theta}}_N$ the Gritsenko--Cl\'ery theta-multiplier.
The $N=6$ case is a theorem (Gritsenko--Nikulin 1998; Cl\'ery);
rows $N\in\{5,7,8\}$ remain open pending a direct Hilbert-series
expansion of $\bigoplus_k M_k(\Gamma^{(2)}_N,\nu^{\mathrm{theta}}_N)$
(Cl\'ery--Gritsenko 2008, 2010).
\end{conjecture}

\begin{center}\small
\begin{tabular}{c|c|c|c|c|c|c}
\hline
$N$ & $(k_N,t_N)$ & Form $\Phi_N$ & Group $\Gamma$ & Multiplier $\nu$ &
$\dim$ & Status \\
\hline
$1$ & $(5,1)$ & $\Delta_5$ & $\Sp_4(\Z)$ & $\nu_{\Delta_5}$ & $1$ &
Theorem (GN 95/98) \\
$2$ & $(2,1)$ & $\Delta_2$ & $\Gamma^{(2)}_2$ & $v_2$ & $1$ &
Theorem (GN 98)\\
$3$ & $(1,1)$ & $\Delta_1$ & $\Gamma^{(2)}_3$ & $v_3$ & $1$ &
Theorem (GN 98)\\
$4$ & $(1,1)$ & $\Delta_{1/2}^{(4)}$ & $\Gamma^{(2)}_4$ & $v_4$ & $1$ &
Theorem (GN 98)\\
$5$ & $(1/2,1)$ & $\Delta_{1/2}^{(5)}$ & $\Gamma^{(2)}_5$ &
$v_5^{\mathrm{theta}}$ & $1$ (conj.) &
Conjecture (CG 08)\\
$6$ & $(1/2,1)$ & $\Delta_{1/2}^{(6)}$ & $\Gamma^{(2)}_6$ &
$v_6^{\mathrm{theta}}$ & $1$ &
Theorem (GN 98; Cl\'ery)\\
$7$ & $(1/2,1)$ & $\Delta_{1/2}^{(7)}$ & $\Gamma^{(2)}_7$ &
$v_7^{\mathrm{theta}}$ & $1$ (conj.) &
Conjecture (Cl\'ery)\\
$8$ & $(1/2,1)$ & $\Delta_{1/2}^{(8)}$ & $\Gamma^{(2)}_8$ &
$v_8^{\mathrm{theta}}$ & $1$ (conj.) &
Conjecture (CG 10)\\
\hline
\end{tabular}
\end{center}

\medskip\noindent
Rows $N\in\{1,2,3,4,6\}$ are exactly those appearing in Borcherds's
$\kBKM(\Phi_N)=c_N(0)/2$ weight table (\emph{Borcherds 1995};
\emph{Gritsenko 1999}, \emph{Duke} 99); rows $N\in\{5,7,8\}$ are the
half-integer-weight extension (Cl\'ery--Gritsenko 2008,
\emph{Acta Arith.}\ 133; 2010, \emph{Nagoya Math.\ J.}\ 199), tying
$M_{1/2}(\Gamma_4,\nu_8)$ and $M_{3/2}(\Gamma_1(4),\nu_4)$ to their
paramodular uplifts.

\subsection{Rigidity and the coincidence of four $c=-214$ derivations}
\label{subsec:c-minus-214-coincidence}

\begin{proposition}[Four $c=-214$ derivations]
\label{prop:four-c-minus-214}
\ClaimStatusTheorem\ (N=1 row only) The four independent chains
\begin{enumerate}[label=\textup{(\roman*)}]
\item Arthur CAP-projection $\pi_{\Delta_5}\hookrightarrow
A_{\mathrm{cusp}}(\Sp_4)$,
\item Kudla--Millson theta lift of the Heegner-divisor class of
$H_1\subset\mathcal{A}_2$,
\item Harvey--Moore heterotic $c=-214$ elliptic genus denominator,
\item Borcherds $c_1(0)/2$ additive lift of $\phi_{0,1}$,
\end{enumerate}
all land in $S_5(\Sp_4(\Z),\nu_{\Delta_5})$. By
Theorem~\ref{thm:dim-S5-one} they equal $C_i\Delta_5$ for scalars
$C_i\in\C^\times$. Each $C_i$ is pinned by the normalisation
$a(T_0)=1$; coincidence $C_1=C_2=C_3=C_4=1$ is a theorem (not a free
match), \emph{provided} the normalisations are tied to the same
Fourier-coefficient datum and the CAP/endoscopic multiplicity is $+1$
(Ta\"ibi 2019).
\end{proposition}

\subsection{Scope declarations}
\label{subsec:scope-rigidity}

\begin{enumerate}
\item The $N=1$ one-dimensionality is a theorem; it rigidifies
$\Delta_5$ up to scalar, and the scalar is fixed by the
Fourier-coefficient normalisation. The $\kBKM(\Phi_1)=5$ identity and
the $\kch(\mathrm{K}3\times E)=\mathrm{wt}(\Delta_5)=5$ matching are
therefore rigidity-forced once $a(T_0)=1$ is assumed.
\item The half-integer-weight rows $N\in\{5,7,8\}$ require
Cl\'ery--Gritsenko theta-multiplier dimension formulas whose published
cases cover $N=6$ fully; $N=5,7,8$ status is
\ClaimStatusConjectured\ until a direct Hilbert-series expansion of
$\bigoplus_k M_k(\Gamma^{(2)}_N,\nu^{\mathrm{theta}}_N)$ is completed.
\item The eightfold Conjecture 1 of Lorgat 2020
(arXiv:2004.01676) requires an eight-tuple of independent arithmetic
rigidities: theorem-status for $N\in\{1,2,3,4,6\}$;
conjecture-status for $N\in\{5,7,8\}$. The unified statement is
therefore \ClaimStatusConjectured\ at present, pending the three
half-weight dimension formulas.
\item The scalar $C_i$ is \emph{not} a free parameter at fixed
Fourier-coefficient normalisation; it is fixed by arithmetic and its
equality across the four chains is a rigidity theorem, not a tautology
of one-dimensionality.
\end{enumerate}

\paragraph{Primary literature.}
Gritsenko--Nikulin, \emph{Igusa modular forms and "the simplest"
Lorentzian Kac--Moody algebras}, \texttt{alg-geom/9504006} (1995);
\emph{Amer.\ J.\ Math.}\ 119 (1997); \emph{Internat.\ J.\ Math.}\ 9
(1998). Borcherds, \emph{Automorphic forms on $O_{s+2,2}(\R)$ and
infinite products}, \emph{Invent.\ Math.}\ 120 (1995). Gritsenko,
\emph{Elliptic genus of Calabi--Yau manifolds and Jacobi and Siegel
modular forms}, \emph{Algebra i Analiz} 11 (1999). Eichler--Zagier,
\emph{The Theory of Jacobi Forms}, Progr.\ Math.\ 55 (1985).
Saito--Kurokawa, \emph{Invent.\ Math.}\ 49 (1978). Pitale--Schmidt,
\emph{JNT} 134 (2014). Weissauer, \emph{Endoscopy for GSp(4) and the
cohomology of Siegel modular threefolds}, LNM 1968 (2009).
Cl\'ery--Gritsenko, \emph{Acta Arith.}\ 133 (2008); \emph{Nagoya Math.\
J.}\ 199 (2010). Lorgat, \emph{arXiv:2004.01676} (2020).
